% Figura: problemas 4.11 y 4.12. Perfiles de campo y potencial frente a la altura.
\begin{figure}[t]
\centering
\begin{subfigure}[b]{0.48\textwidth}\centering
\begin{tikzpicture}
\begin{axis}[width=6.6cm, height=5.4cm, xmin=0, xmax=1.65, ymin=0, ymax=2,
  ytick={0,1,2}, yticklabels={$0$,$a$,$d$}, ylabel={altura}, xtick={0,0.5,1,1.5},
  xticklabels={$0$,$\frac12$,$1$,$\frac32$}, grid=major,
  legend style={at={(0.5,1.03)}, anchor=south, legend columns=3, font=\tiny, /tikz/every even column/.append style={column sep=3pt}}]
  \fill[dielec, opacity=0.6] (axis cs:0,0) rectangle (axis cs:1.65,1);
  \addplot[curvaD, line width=2.8pt, opacity=0.45] coordinates {(1.5,0) (1.5,2)}; \addlegendentry{$D\,d/(\eo V_0)$}
  \addplot[curvaE] coordinates {(0.5,0) (0.5,1)}; \addlegendentry{$E\,d/V_0$}
  \addplot[curvaV] coordinates {(0,0) (0.25,1) (1,2)}; \addlegendentry{$\phi/V_0$}
  \addplot[curvaE, forget plot] coordinates {(1.5,1) (1.5,2)};
  \addplot[curvaE, forget plot, dotted, line width=0.5pt] coordinates {(0.5,1) (1.5,1)};
  \node[etiqt, anchor=north east] at (axis cs:1.35,0.95) {$\varepsilon=3\eo$};
\end{axis}
\end{tikzpicture}
\caption{Problema 4.11 ($a=b$, $\er=3$)}
\end{subfigure}\hfill
\begin{subfigure}[b]{0.48\textwidth}\centering
\begin{tikzpicture}
\begin{axis}[width=6.6cm, height=5.4cm, xmin=0, xmax=1.4, ymin=0, ymax=3,
  ytick={0,1,1.6,3}, yticklabels={$0$,$a$,$a+e$,$d$}, ylabel={altura}, xtick={0,0.5,1},
  xticklabels={$0$,$\frac12$,$1$}, grid=major,
  legend style={at={(0.5,1.03)}, anchor=south, legend columns=3, font=\tiny, /tikz/every even column/.append style={column sep=3pt}}]
  \fill[conductor, opacity=0.45] (axis cs:0,1) rectangle (axis cs:1.4,1.6);
  \addplot[curvaE] coordinates {(1,0) (1,1) (0,1) (0,1.6) (1,1.6) (1,3)}; \addlegendentry{$E\,(d-e)/V_0$}
  \addplot[curvaV] coordinates {(0,0) (0.4167,1) (0.4167,1.6) (1,3)}; \addlegendentry{$\phi/V_0$}
  \node[etiqt] at (axis cs:0.95,1.3) {metal};
\end{axis}
\end{tikzpicture}
\caption{Problema 4.12 ($a=1$, $e=0{,}6$, $b=1{,}4$)}
\end{subfigure}
\caption{Perfiles en función de la altura sobre la placa inferior (que se toma a potencial nulo). (a) Capa dieléctrica de
espesor $a$ bajo una capa de aire de espesor $b$: $D$ es continuo, el campo salta en el factor $\er$ y el potencial
cae sobre todo en el aire. (b) Lámina conductora: el campo es nulo dentro del metal y el potencial es allí
constante; fuera, el campo es el de un condensador de separación $d-e$.}
\label{fig:p411-perfil}
\end{figure}
